\chapter{Proof of Theorem 1}
\label{proof_theorem_1}
The appendix of \citeA{hattSequentialDeconfoundingCausal2024} 
formulates the proof for Theorem 1 only for a first-order
conditional Markov model over treatments. 
However, the same proof strategy extends naturally
to finite order \(p\) by augmenting the treatment state with the 
last \(p\) treatments. The
result below is therefore a straightforward extension
of their recursive-construction argument. 

\begin{proposition}[Finite-order extension of the proof strategy of \citeA{hattSequentialDeconfoundingCausal2024}]
Suppose that for some fixed integer \(p \ge 1\), the assigned treatment process satisfies
the conditional Markov property of order \(p\), that is,
\begin{equation}
\label{eq:p-order-cmm}
p_{\boldsymbol{\gamma}}(\bar{\mathbf{a}}_{iT}\mid \bar{\mathbf{x}}_{iT},\mathbf{z}_i)
=
\prod_{t=1}^T
p_{\boldsymbol{\gamma}_t}(a_{it}\mid \mathbf{a}_{i,t-p:t-1},\bar{\mathbf{x}}_{it},\mathbf{z}_i),
\end{equation}
where \(\mathbf{a}_{i,t-p:t-1}=(a_{i,t-p},\dots,a_{i,t-1})\), with the convention that for \(t<1\),
\(a_{it}=0\). Assume further that the treatment space is a Borel space and that unobserved
confounding is time-invariant in the sense maintained by \citeA{hattSequentialDeconfoundingCausal2024}.

Then the proof strategy of \citeA{hattSequentialDeconfoundingCausal2024} extends to yield a recursive representation
\begin{equation}
\label{eq:p-order-recursion}
A_{it} = f_t(\mathbf{A}_{i,t-p:t-1},\mathbf{Z}_i,\bar{\mathbf{X}}_{it},V_{it}),
\end{equation}
where \(V_{it} \sim \mathrm{Uniform}([0,1])\), and
\begin{equation}
\label{eq:p-order-v-indep}
V_{it} \perp Y_i(\bar{\mathbf{a}}_t)\mid \bar{\mathbf{A}}_{i,t-1},\bar{\mathbf{X}}_{it},\mathbf{Z}_i
\qquad
\text{for all } \bar{\mathbf{a}}_t \in \bar{\mathcal A}_t.
\end{equation}
Consequently,
\begin{equation}
\label{eq:p-order-ssi}
A_{it} \perp Y_i(\bar{\mathbf{a}}_t)\mid \bar{\mathbf{A}}_{i,t-1},\bar{\mathbf{X}}_{it},\mathbf{Z}_i
\qquad
\text{for all } \bar{\mathbf{a}}_t \in \bar{\mathcal A}_t.
\end{equation}
\end{proposition}

\begin{proof}
The argument follows the same structure as the appendix of \citeA{hattSequentialDeconfoundingCausal2024}, but
with a state augmentation step.

First, define the augmented treatment state
\[
\mathbf{S}_{it} := (A_{i,t-p+1},\dots, A_{it}) \in \mathcal A^p.
\]
Under \eqref{eq:p-order-cmm}, the treatment process is Markov of order \(p\) conditional on
\((\bar{\mathbf{X}}_{it},\mathbf{Z}_i)\). Equivalently, the augmented state process \((\mathbf{S}_{it})_{t\ge1}\) is first-order
Markov conditional on \((\bar{\mathbf{X}}_{it},\mathbf{Z}_i)\). Indeed, for each \(t\),
\[
\mathbb{P}(\mathbf{S}_{it} \in B \mid \bar{\mathbf{S}}_{i,t-1}, \bar{\mathbf{X}}_{it}, \mathbf{Z}_i)
=
\mathbb{P}(\mathbf{S}_{it} \in B \mid \mathbf{S}_{i,t-1}, \bar{\mathbf{X}}_{it}, \mathbf{Z}_i)
\]
for all measurable \(B \subseteq \mathcal A^p\), because the first \(p-1\) coordinates of
\(\mathbf{S}_{it}\) are deterministic shifts of \(\mathbf{S}_{i,t-1}\), while the last coordinate \(A_{it}\) depends
on the past only through \(\mathbf{A}_{i,t-p:t-1}\), \(\bar{\mathbf{X}}_{it}\), and \(\mathbf{Z}_i\).

Since \(\mathcal A^p\) is again a Borel space, the same Kallenberg recursion argument used by
\citeA{hattSequentialDeconfoundingCausal2024} for the order-1 case applies to the augmented state process.
Therefore, there exist measurable functions \(F_t\) and i.i.d. random variables
\(V_{it} \sim \mathrm{Uniform}([0,1])\) such that
\[
\mathbf{S}_{it} = F_t(\mathbf{S}_{i,t-1}, \mathbf{Z}_i, \bar{\mathbf{X}}_{it}, V_{it}).
\]
Let \(f_t\) denote the last-coordinate projection of \(F_t\). Then
\[
A_{it} = f_t(\mathbf{A}_{i,t-p:t-1}, \mathbf{Z}_i, \bar{\mathbf{X}}_{it}, V_{it}),
\]
which establishes \eqref{eq:p-order-recursion}.

It remains to verify \eqref{eq:p-order-v-indep}. This is the same step that remains to be shown
in the proof of Lemma 2 in \citeA{hattSequentialDeconfoundingCausal2024}. Their argument extends without essential
change. Suppose, for contradiction, that there exists an omitted variable $W$ that confounds
\(V_{it}\) and \(Y_i(\bar{\mathbf{a}}_t)\) conditional on \((\bar{\mathbf{A}}_{i,t-1},\bar{\mathbf{X}}_{it},\mathbf{Z}_i)\). 
Moreover, according to Assumption~\ref{ass:SG-No-single-time-point-confounding}, there is no 
single-time confounders. Hence, $W$ at least affect $V_{is}$ where $s \neq t$. 
But the recursion construction yields that the \(V_{it}\) are i.i.d. Uniform\([0,1]\),
so such a common time-invariant source of dependence is excluded. 
If instead the omitted variable is time-varying, then because
\[
A_{it} = f_t(\mathbf{A}_{i,t-p:t-1}, \mathbf{Z}_i, \bar{\mathbf{X}}_{it}, V_{it}),
\]
the same omitted variable would induce confounding between \(A_{it}\) and \(Y_i(\bar{\mathbf{a}}_t)\).
This contradicts the maintained setup that any unobserved confounding is time-invariant and is
captured through the latent static variable \(\mathbf{Z}_i\). Hence,
\[
V_{it} \perp Y_i(\bar{\mathbf{a}}_t)\mid \bar{\mathbf{A}}_{i,t-1},\bar{\mathbf{X}}_{it},\mathbf{Z}_i,
\]
which proves \eqref{eq:p-order-v-indep}.

Finally, define
\[
\mathcal G_{it} := \sigma(\bar{\mathbf{A}}_{i,t-1},\bar{\mathbf{X}}_{it},\mathbf{Z}_i).
\]
Since \((\mathbf{A}_{i,t-p:t-1}, \mathbf{Z}_i, \bar{\mathbf{X}}_{it})\) is \(\mathcal G_{it}\)-measurable, \eqref{eq:p-order-v-indep}
implies
\[
(\mathbf{A}_{i,t-p:t-1}, \mathbf{Z}_i, \bar{\mathbf{X}}_{it}, V_{it}) \perp Y_i(\bar{\mathbf{a}}_t)\mid \mathcal G_{it}.
\]
Because \(A_{it}\) is a measurable function of \((\mathbf{A}_{i,t-p:t-1}, \mathbf{Z}_i, \bar{\mathbf{X}}_{it}, V_{it})\), conditional
independence is preserved under measurable transformations, and therefore
\[
A_{it} \perp Y_i(\bar{\mathbf{a}}_t)\mid \mathcal G_{it}.
\]
Equivalently,
\[
A_{it} \perp Y_i(\bar{\mathbf{a}}_t)\mid \bar{\mathbf{A}}_{i,t-1},\bar{\mathbf{X}}_{it},\mathbf{Z}_i.
\]
This proves \eqref{eq:p-order-ssi}.
\end{proof}
